% Options for packages loaded elsewhere
\PassOptionsToPackage{unicode}{hyperref}
\PassOptionsToPackage{hyphens}{url}
\documentclass[
]{article}
\usepackage{xcolor}
\usepackage[margin=1in]{geometry}
\usepackage{amsmath,amssymb}
\setcounter{secnumdepth}{-\maxdimen} % remove section numbering
\usepackage{iftex}
\ifPDFTeX
  \usepackage[T1]{fontenc}
  \usepackage[utf8]{inputenc}
  \usepackage{textcomp} % provide euro and other symbols
\else % if luatex or xetex
  \usepackage{unicode-math} % this also loads fontspec
  \defaultfontfeatures{Scale=MatchLowercase}
  \defaultfontfeatures[\rmfamily]{Ligatures=TeX,Scale=1}
\fi
\usepackage{lmodern}
\ifPDFTeX\else
  % xetex/luatex font selection
\fi
% Use upquote if available, for straight quotes in verbatim environments
\IfFileExists{upquote.sty}{\usepackage{upquote}}{}
\IfFileExists{microtype.sty}{% use microtype if available
  \usepackage[]{microtype}
  \UseMicrotypeSet[protrusion]{basicmath} % disable protrusion for tt fonts
}{}
\makeatletter
\@ifundefined{KOMAClassName}{% if non-KOMA class
  \IfFileExists{parskip.sty}{%
    \usepackage{parskip}
  }{% else
    \setlength{\parindent}{0pt}
    \setlength{\parskip}{6pt plus 2pt minus 1pt}}
}{% if KOMA class
  \KOMAoptions{parskip=half}}
\makeatother
\usepackage{color}
\usepackage{fancyvrb}
\newcommand{\VerbBar}{|}
\newcommand{\VERB}{\Verb[commandchars=\\\{\}]}
\DefineVerbatimEnvironment{Highlighting}{Verbatim}{commandchars=\\\{\}}
% Add ',fontsize=\small' for more characters per line
\usepackage{framed}
\definecolor{shadecolor}{RGB}{248,248,248}
\newenvironment{Shaded}{\begin{snugshade}}{\end{snugshade}}
\newcommand{\AlertTok}[1]{\textcolor[rgb]{0.94,0.16,0.16}{#1}}
\newcommand{\AnnotationTok}[1]{\textcolor[rgb]{0.56,0.35,0.01}{\textbf{\textit{#1}}}}
\newcommand{\AttributeTok}[1]{\textcolor[rgb]{0.13,0.29,0.53}{#1}}
\newcommand{\BaseNTok}[1]{\textcolor[rgb]{0.00,0.00,0.81}{#1}}
\newcommand{\BuiltInTok}[1]{#1}
\newcommand{\CharTok}[1]{\textcolor[rgb]{0.31,0.60,0.02}{#1}}
\newcommand{\CommentTok}[1]{\textcolor[rgb]{0.56,0.35,0.01}{\textit{#1}}}
\newcommand{\CommentVarTok}[1]{\textcolor[rgb]{0.56,0.35,0.01}{\textbf{\textit{#1}}}}
\newcommand{\ConstantTok}[1]{\textcolor[rgb]{0.56,0.35,0.01}{#1}}
\newcommand{\ControlFlowTok}[1]{\textcolor[rgb]{0.13,0.29,0.53}{\textbf{#1}}}
\newcommand{\DataTypeTok}[1]{\textcolor[rgb]{0.13,0.29,0.53}{#1}}
\newcommand{\DecValTok}[1]{\textcolor[rgb]{0.00,0.00,0.81}{#1}}
\newcommand{\DocumentationTok}[1]{\textcolor[rgb]{0.56,0.35,0.01}{\textbf{\textit{#1}}}}
\newcommand{\ErrorTok}[1]{\textcolor[rgb]{0.64,0.00,0.00}{\textbf{#1}}}
\newcommand{\ExtensionTok}[1]{#1}
\newcommand{\FloatTok}[1]{\textcolor[rgb]{0.00,0.00,0.81}{#1}}
\newcommand{\FunctionTok}[1]{\textcolor[rgb]{0.13,0.29,0.53}{\textbf{#1}}}
\newcommand{\ImportTok}[1]{#1}
\newcommand{\InformationTok}[1]{\textcolor[rgb]{0.56,0.35,0.01}{\textbf{\textit{#1}}}}
\newcommand{\KeywordTok}[1]{\textcolor[rgb]{0.13,0.29,0.53}{\textbf{#1}}}
\newcommand{\NormalTok}[1]{#1}
\newcommand{\OperatorTok}[1]{\textcolor[rgb]{0.81,0.36,0.00}{\textbf{#1}}}
\newcommand{\OtherTok}[1]{\textcolor[rgb]{0.56,0.35,0.01}{#1}}
\newcommand{\PreprocessorTok}[1]{\textcolor[rgb]{0.56,0.35,0.01}{\textit{#1}}}
\newcommand{\RegionMarkerTok}[1]{#1}
\newcommand{\SpecialCharTok}[1]{\textcolor[rgb]{0.81,0.36,0.00}{\textbf{#1}}}
\newcommand{\SpecialStringTok}[1]{\textcolor[rgb]{0.31,0.60,0.02}{#1}}
\newcommand{\StringTok}[1]{\textcolor[rgb]{0.31,0.60,0.02}{#1}}
\newcommand{\VariableTok}[1]{\textcolor[rgb]{0.00,0.00,0.00}{#1}}
\newcommand{\VerbatimStringTok}[1]{\textcolor[rgb]{0.31,0.60,0.02}{#1}}
\newcommand{\WarningTok}[1]{\textcolor[rgb]{0.56,0.35,0.01}{\textbf{\textit{#1}}}}
\usepackage{graphicx}
\makeatletter
\newsavebox\pandoc@box
\newcommand*\pandocbounded[1]{% scales image to fit in text height/width
  \sbox\pandoc@box{#1}%
  \Gscale@div\@tempa{\textheight}{\dimexpr\ht\pandoc@box+\dp\pandoc@box\relax}%
  \Gscale@div\@tempb{\linewidth}{\wd\pandoc@box}%
  \ifdim\@tempb\p@<\@tempa\p@\let\@tempa\@tempb\fi% select the smaller of both
  \ifdim\@tempa\p@<\p@\scalebox{\@tempa}{\usebox\pandoc@box}%
  \else\usebox{\pandoc@box}%
  \fi%
}
% Set default figure placement to htbp
\def\fps@figure{htbp}
\makeatother
\setlength{\emergencystretch}{3em} % prevent overfull lines
\providecommand{\tightlist}{%
  \setlength{\itemsep}{0pt}\setlength{\parskip}{0pt}}
\usepackage{mathpazo}
\usepackage{booktabs}
\usepackage{rotating}
\usepackage{makecell}
\usepackage{hyperref}
%% \usepackage[a4paper,margin=2cm]{geometry}


\newcommand{\orcidlogo}{\textcolor{orcidlogocol}{\aiOrcid}}
\usepackage[font=Large,labelfont=bf,textfont=bf]{caption}
\DeclareCaptionLabelFormat{addC}{#1 C#2}


%% This solution for adding the ORCID logo was taken from:
%% https://tex.stackexchange.com/questions/445563/ieeetran-how-to-include-orcid-in-tex-pdf-with-pdflatex/445583
%% and adapted simply by scaling the icon a little smaller than | to +
\usepackage{scalerel}
\usepackage{tikz}
\usetikzlibrary{svg.path}

\definecolor{orcidlogocol}{HTML}{A6CE39}
\tikzset{
  orcidlogo/.pic={
    \fill[orcidlogocol] svg{M256,128c0,70.7-57.3,128-128,128C57.3,256,0,198.7,0,128C0,57.3,57.3,0,128,0C198.7,0,256,57.3,256,128z};
    \fill[white] svg{M86.3,186.2H70.9V79.1h15.4v48.4V186.2z}
                 svg{M108.9,79.1h41.6c39.6,0,57,28.3,57,53.6c0,27.5-21.5,53.6-56.8,53.6h-41.8V79.1z M124.3,172.4h24.5c34.9,0,42.9-26.5,42.9-39.7c0-21.5-13.7-39.7-43.7-39.7h-23.7V172.4z}
                 svg{M88.7,56.8c0,5.5-4.5,10.1-10.1,10.1c-5.6,0-10.1-4.6-10.1-10.1c0-5.6,4.5-10.1,10.1-10.1C84.2,46.7,88.7,51.3,88.7,56.8z};
  }
}

\newcommand\orcidicon[1]{\href{https://orcid.org/#1}{\mbox{\scalerel*{
\begin{tikzpicture}[yscale=-1,transform shape]
\pic{orcidlogo};
\end{tikzpicture}
}{+}}}}

%% Set hyperlinks in colour to the "CODECHECK colour".
\hypersetup{colorlinks=true, urlcolor = [rgb]{0,.5,0.2}}
\usepackage{bookmark}
\IfFileExists{xurl.sty}{\usepackage{xurl}}{} % add URL line breaks if available
\urlstyle{same}
\hypersetup{
  pdftitle={CODECHECK certificate 2020-XYZ},
  hidelinks,
  pdfcreator={LaTeX via pandoc}}

\title{CODECHECK certificate 2020-XYZ}
\usepackage{etoolbox}
\makeatletter
\providecommand{\subtitle}[1]{% add subtitle to \maketitle
  \apptocmd{\@title}{\par {\large #1 \par}}{}{}
}
\makeatother
\subtitle{\url{https://doi.org/10.5281/zenodo.FIXME}}
\author{}
\date{\vspace{-2.5em}}

\begin{document}
\maketitle

\centerline{\includegraphics[width=4cm]{/Library/Frameworks/R.framework/Versions/4.4-arm64/Resources/library/codecheck/extdata/codecheck_logo.pdf}}\vspace*{2cm}

\begin{table}[ht]
\centering
\begin{tabular}{lp{10cm}}
  \hline
Item & Value \\ 
  \hline
Title & Quality in Question: Assessing the accuracy of four heart rate wearables and the implications for psychophysiological research.
 \\ 
  Authors & Mohammadamin Sinichi \orcidicon{0009-0008-2491-1542} , Martin Gevonden, Lydia Krabbendam \\ 
  Reference & \url{https://doi.org/10.31219/osf.io/wkzsn} \\ 
  Codechecker & Yasel Quintero \orcidicon{0009-0005-9240-7370} , Tornike Skhirtladze, Joslyn Sun \orcidicon{0009-0002-2379-0060} , Gabriella Low Chew Tung, Roel Janssen \orcidicon{0000-0003-4324-5350}  \\ 
  Date of check & 2020-11-28 14:00:00 \\ 
  Summary & Jupyter notebooks, Python package and data were provided by the author. We edited the paths to the .CSV data files in the notebooks to reproduce Figure 5 in the manuscript.
 \\ 
  Repository & \url{https://github.com/yiquintero/wearable-hrv} \\ 
   \hline
\end{tabular}
\caption{CODECHECK summary} 
\end{table}

\begin{table}[ht]
\centering
\begin{tabular}{p{6cm}p{6cm}p{2cm}}
  \hline
Output & Comment & Size (b) \\ 
  \hline
\href{https://github.com/yiquintero/wearable-hrv/blob/master/codecheck/outputs/Figure5.png}{\path{figures-codecheck/Figure5.png}} & manuscript Figure 5 & 337355 \\ 
   \hline
\end{tabular}
\caption{Summary of output files generated} 
\end{table}

\section{Summary}\label{summary}

This code was (not) straightforward to codecheck. {[}\ldots{} ADD MORE
INTERESTING FINDINGS HERE..{]}

\clearpage

\section{CODECHECKER notes}\label{codechecker-notes}

This took \ldots{} minutes to complete on \{a large workstation, my
laptop\}.

\subsection{Recommendations}\label{recommendations}

I suggest to the authors to consider the following suggestions for their
next publication or workflow:

\begin{itemize}
\tightlist
\item
  \ldots{}
\end{itemize}

(document here if any of the suggestions were taken up by the authors in
the meantime - do not remove any, keep track of contributions via
feedback)

\clearpage

\section{Manifest files}\label{manifest-files}

\subsection{Figure5.png}\label{figure5.png}

\textbf{Comment:} manuscript Figure 5

\pandocbounded{\includegraphics[keepaspectratio]{/Users/yquinterolares/repos/codecheck/wearable-hrv/codecheck/outputs/Figure5.png}}
\clearpage 

\clearpage

\subsection{Acknowledgements}\label{acknowledgements}

I would like to thank Dr Bhatt and his team for promptly answering any
queries I had with this reproduction. CODECHECK is financially supported
by the Mozilla foundation.

\section{Citing this document}\label{citing-this-document}

Yasel Quintero, Tornike Skhirtladze, Joslyn Sun, Gabriella Low Chew
Tung, Roel Janssen (2020). CODECHECK Certificate 2020-XYZ. Zenodo.
\url{https://doi.org/10.5281/zenodo.FIXME}

\section{About CODECHECK}\label{about-codecheck}

This certificate confirms that the codechecker could independently
reproduce the results of a computational analysis given the data and
code from a third party. A CODECHECK does not check whether the original
computation analysis is correct. However, as all materials required for
the reproduction are freely available by following the links in this
document, the reader can then study for themselves the code and data.

\section{About this document}\label{about-this-document}

This document was created using \href{https://rmarkdown.rstudio.com/}{R
Markdown} using the
\href{https://github.com/codecheckers/codecheck}{\texttt{codecheck}} R
package. \texttt{make\ codecheck.pdf} will regenerate the report file.

\begin{Shaded}
\begin{Highlighting}[]
\FunctionTok{sessionInfo}\NormalTok{()}
\end{Highlighting}
\end{Shaded}

\begin{verbatim}
## R version 4.4.1 (2024-06-14)
## Platform: aarch64-apple-darwin20
## Running under: macOS 15.1.1
## 
## Matrix products: default
## BLAS:   /Library/Frameworks/R.framework/Versions/4.4-arm64/Resources/lib/libRblas.0.dylib 
## LAPACK: /Library/Frameworks/R.framework/Versions/4.4-arm64/Resources/lib/libRlapack.dylib;  LAPACK version 3.12.0
## 
## locale:
## [1] C/UTF-8/C/C/C/C
## 
## time zone: Europe/Amsterdam
## tzcode source: internal
## 
## attached base packages:
## [1] stats     graphics  grDevices utils     datasets 
## [6] methods   base     
## 
## other attached packages:
##  [1] readr_2.1.5      tibble_3.2.1     xtable_1.8-4    
##  [4] yaml_2.3.10      rprojroot_2.0.4  knitr_1.49      
##  [7] codecheck_0.11.4 parsedate_1.3.1  R.cache_0.16.0  
## [10] gh_1.4.1        
## 
## loaded via a namespace (and not attached):
##  [1] xfun_0.49         rdflib_0.2.9      tzdb_0.4.0       
##  [4] vctrs_0.6.5       tools_4.4.1       generics_0.1.3   
##  [7] curl_6.0.1        parallel_4.4.1    fansi_1.0.6      
## [10] pkgconfig_2.0.3   pdftools_3.4.1    R.oo_1.27.0      
## [13] redland_1.0.17-18 assertthat_0.2.1  lifecycle_1.0.4  
## [16] compiler_4.4.1    atom4R_0.3-3      stringr_1.5.1    
## [19] keyring_1.3.2     htmltools_0.5.8.1 pillar_1.9.0     
## [22] crayon_1.5.3      whisker_0.4.1     tidyr_1.3.1      
## [25] R.utils_2.12.3    cachem_1.1.0      zen4R_0.10       
## [28] tidyselect_1.2.1  zip_2.3.1         digest_0.6.37    
## [31] stringi_1.8.4     dplyr_1.1.4       purrr_1.0.2      
## [34] fastmap_1.2.0     cli_3.6.3         magrittr_2.0.3   
## [37] XML_3.99-0.17     crul_1.5.0        utf8_1.2.4       
## [40] osfr_0.2.9        withr_3.0.2       bit64_4.5.2      
## [43] roxygen2_7.3.2    rmarkdown_2.29    httr_1.4.7       
## [46] bit_4.5.0         qpdf_1.3.4        askpass_1.2.1    
## [49] R.methodsS3_1.8.2 hms_1.1.3         memoise_2.0.1    
## [52] evaluate_1.0.1    rlang_1.1.4       Rcpp_1.0.13-1    
## [55] glue_1.8.0        httpcode_0.3.0    xml2_1.3.6       
## [58] fauxpas_0.5.2     rorcid_0.7.0      vroom_1.6.5      
## [61] jsonlite_1.8.9    plyr_1.8.9        R6_2.5.1         
## [64] fs_1.6.5
\end{verbatim}

\end{document}
